\input{../tex-inputs/latex-leseansicht-vorspann}

\section[Gustav Schwarzkopf an Arthur Schnitzler, 9. 8. 1899]{L04296 Gustav Schwarzkopf an Arthur Schnitzler, 9. 8. 1899}
\nopagebreak\mylabel{L04296v}
\rehead{ }\normalsize\beginnumbering\briefempfaengerindex{Schnitzler, Arthur@\textsc{Schnitzler, Arthur}!zzzSchwarzkopf, Gustav@\emph{von Gustav Schwarzkopf}!1899-08-091@{9. 8. 1899}|(be}
\toendnotes[C]{\smallbreak\pagebreak[2]}
\correspDesc{Versand  durch Gustav Schwarzkopf am 9. 8. 1899 in Wien
\newline{}Übermittlung  am 9. 8. 1899 in Wien
\newline{}Zustellung  am 11. 8. 1899 in Bozen
\newline{}Erhalt  durch Arthur Schnitzler am 12. 8. 1899 in Bozen}\toendnotes[C]{\smallbreak}
\Standort{CUL, Schnitzler, B 96a.}
\physDesc{Postkarte, 874 Zeichen
\newline{}Handschrift: schwarze Tinte, deutsche Kurrent
\newline{}Versand: 1) Stempel: »\nobreak{}\oindex{I., Innere Stadt@\textbf{I., Innere Stadt}, \emph{Verwaltungsgebiet}|pwk}Wien 1/1 1, 9. 8. 99, 4N\nobreak{}«.   2) Stempel: »\nobreak{}\oindex{Bozen@\textbf{Bozen}, \emph{Hauptstadt}|pwk}B{[}oz{]}\textcolor{gray}{en}
                                       1, 11. 8. 99, 7.F\nobreak{}«. 
\newline{}Schnitzler: mit Bleistift datiert:
                                       »9. 8. 9\textcolor{gray}{9}« }\toendnotes[C]{\smallbreak}\pstart{}{\pb}Herrn \textsc{D\textsuperscript{r} Arthur Schnitzler}\pend{}\pstart{}\textsc{Bozen\oindex{Bozen@\textbf{Bozen}, \emph{Hauptstadt}|pw}}\pend{}\pstart{}\textsc{poste. rest.}\pend{}{\bigskip}\vspace{1em}
\pstart
           \noindent{}{\pb}Lieber Arthur, ich hatte urſprünglich die Abſicht, die in die
                  \label{K_L04296-1v}\edtext{Civiliſation zurückkehrenden kühnen
                  Forſcher\pwindex{Beer-Hofmann, Richard 11.\,7.\,1866 Wien – 26.\,9.\,1945 New York City@\textsc{Beer-Hofmann, Richard} (11.\,7.\,1866 Wien – 26.\,9.\,1945 New York City), \emph{Schriftsteller}|pwv}\pwindex{Wassermann, Jakob 10.\,3.\,1873 Fürth – 1.\,1.\,1934 Altaussee@\textsc{Wassermann, Jakob} (10.\,3.\,1873 Fürth – 1.\,1.\,1934 Altaussee), \emph{Schriftsteller}|pwv}}{\lemma{\textnormal{\emph{Civilisation … Forscher}}}\Cendnote{\textnormal{Zwischen 5. 8. 1899 und 5. 8. 1899 unternahm Schnitzler mit Richard
                     Beer-Hofmann\pwindex{Beer-Hofmann, Richard 11.\,7.\,1866 Wien – 26.\,9.\,1945 New York City@\textsc{Beer-Hofmann, Richard} (11.\,7.\,1866 Wien – 26.\,9.\,1945 New York City), \emph{Schriftsteller}|pwk} und Jakob Wassermann\pwindex{Wassermann, Jakob 10.\,3.\,1873 Fürth – 1.\,1.\,1934 Altaussee@\textsc{Wassermann, Jakob} (10.\,3.\,1873 Fürth – 1.\,1.\,1934 Altaussee), \emph{Schriftsteller}|pwk}
                  eine Fußwanderung durch Südtirol\oindex{Südtirol@\textbf{Südtirol}, \emph{Verwaltungsgebiet}|pwk}. }}}\label{K_L04296-1} in
               einem{ }ſchwungvollen Feſtgedicht zu begrüßen, wie das{ }ſo meine Art iſt. Aber es fehlt
               mir das Material, ich weiß nicht, welche Vaſe Sie entdeckt, welche Leiſtungen Sie
               vollbracht haben. Keine Anſichtskarte hat den Weg zu mir gefunden. Und \label{T_L04296-1v}\edtext{Sie}{\lemma{\textnormal{\emph{Sie}}}\Cendnote{\textnormal{Er schreibt: »ſie«.}}}\label{T_L04296-1} waren doch beſti{\geminationm}t den{ }ſtolzen Anfang eines Albums zu bilden! Na, ich
               nehme an, daß alles gut gegangen iſt – die Herren auch – und daß die verſchiedenen
               Damen die beſten Fortſchritte gemacht haben. We{\geminationn} Sie es
               übrigens{ }ſo heiß gehabt haben wie wir in Wien\oindex{Wien@\textbf{Wien}, \emph{Verwaltungsgebiet}|pw},
                  da{\geminationn} muß es nicht i{\geminationm}er
               ein ungetrübtes Vergnügen geweſen{ }ſein.– Ich war in der vorigen Woche 2 ½ Tage
               in der Brühl\oindex{Brühl@\textbf{Brühl}, \emph{Tal}|pw} und hätte von dort beinahe einen
               Ausflug auf den Huſarentempel\oindex{Husarentempel@\textbf{Husarentempel}, \emph{Monument}|pw} gemacht.\pend
           
\pstart
           Herzlichſt{\\[\baselineskip]}Ihr{\\[\baselineskip]}\spacefill\mbox{GustavSch}\pend
           \leftskip=0em{}
\pstart
           \noindent{}Gruß an Richard\pwindex{Beer-Hofmann, Richard 11.\,7.\,1866 Wien – 26.\,9.\,1945 New York City@\textsc{Beer-Hofmann, Richard} (11.\,7.\,1866 Wien – 26.\,9.\,1945 New York City), \emph{Schriftsteller}|pw} und Herr Waſſerma{\geminationn}\pwindex{Wassermann, Jakob 10.\,3.\,1873 Fürth – 1.\,1.\,1934 Altaussee@\textsc{Wassermann, Jakob} (10.\,3.\,1873 Fürth – 1.\,1.\,1934 Altaussee), \emph{Schriftsteller}|pw}\pend
           \selectlanguage{ngerman}\endnumbering\briefempfaengerindex{Schnitzler, Arthur@\textsc{Schnitzler, Arthur}!zzzSchwarzkopf, Gustav@\emph{von Gustav Schwarzkopf}!1899-08-091@{9. 8. 1899}|)be}\mylabel{L04296h}
\begin{anhang}
\end{anhang}\newcommand{\dateiname}{L04296}\newcommand{\titel}{Gustav Schwarzkopf an Arthur Schnitzler, 9. 8. 1899}\newcommand{\editorInnen}{Herausgegeben von Jahnke, SelmaMüller, Martin Anton}\input{../tex-inputs/latex-leseansicht-abspann}
